\BourbakiExerciseGroup

\begin{enumerate}
\def\labelenumi{\arabic{enumi})}
\item
  在豪斯多夫拓扑环 A 中，A 的每个子集的中心化子（特别地，A 的中心）是闭的，并且 A 的每个子集的左（相应地，右）零化子也是闭的。
\item
  \begin{enumerate}
  \def\labelenumii{\alph{enumii})}
  \tightlist
  \item
    设 \(\mathfrak{a}\) 是拓扑环 A 中的离散左理想。试证：对每个 \(x \in \mathfrak{a}\) ，\(x\) 在 A 中的左零化子是开的。由此试证：若 A 是无非零因子的非离散环，则它不含 \(\{\mathbf{o}\}\) 以外的离散（左或右）理想。
  \end{enumerate}
\end{enumerate}

\begin{enumerate}
\def\labelenumi{\alph{enumi})}
\setcounter{enumi}{1}
\tightlist
\item
  设 A 是局部连通拓扑环，\(\mathfrak{a}\) 是 A 的全不连通左理想。试证：对每个 \(x \in \mathfrak{a}\) ，\(x\) 在 A 中的左零化子是开的。由此试证：若 A 无零因子，则它不含 \(\{0\}\) 以外的全不连通（左或右）理想。
\end{enumerate}

\begin{enumerate}
\def\labelenumi{\arabic{enumi})}
\setcounter{enumi}{2}
\item
  拓扑环 A 中 o 的连通分支是双边理想。由此推出：若 A 是拟单的且不连通（特别地，若 A 是不连通的拓扑除环），则 A 全不连通。
\item
  设 \((x_{i})\) 是豪斯多夫拓扑环 A 中的可和族，\(s\) 为其和。则对每个 \(a \in A\) ，族 \((ax_{i})\){[}相应地 \((x_{i}a)]\)可和，且其和为 as（相应地，sa）。若 \((x_{\lambda})\) 与 \((y_{\mu})\) 是 A 中两个可和族，且族 \((x_{\lambda}y_{\mu})\) 可和，则 \(\sum_{\lambda,\mu} x_{\lambda} y_{\mu} = \left( \sum_{\lambda} x_{\lambda} \right) \left( \sum_{\mu} y_{\mu} \right)\) 。此外，若 A 完备且诸 \(x_{\lambda}\) 之一是可逆元，则族 \((x_{\lambda}y_{\mu})\) 的可和性蕴含 \((y_{\mu})\) 的可和性。
\item
  矩形是 \(N \times N\) 的一个子集，它是 N 的两个区间的乘积。给定 \(N \times N\) 的有限子集 E，令 \(\varphi(E)\) 为其并为 E 的两两不相交矩形的最小个数。设 \((E_n)\) 是 \(N \times N\) 的有限子集的一个递增序列，它覆盖 \(N \times N\) ，且使序列 \((\varphi(E_n))\) 有界。若 \((x_n), (y_n)\) 是两个收敛级数，其项都属于豪斯多夫拓扑环 A，则我们有
\end{enumerate}

\[
\lim _ {n \rightarrow \infty} \sum_ {(h, k) \in \mathrm{E} _ {n}} x _ {h} y _ {k} = \left(\underset {n = 0} {\overset {\infty} {\sum}} x _ {n}\right)\left(\underset {n = 0} {\overset {\infty} {\sum}} y _ {n}\right).
\]

试构造一个满足上述条件的序列 \((\mathbf{E}_{n})\) 的例子，使得对每个 n，\(E_{n+1} - E_{n}\) 只含一个元素。

\begin{enumerate}
\def\labelenumi{\arabic{enumi})}
\setcounter{enumi}{5}
\item
  设 A 是有单位元的拓扑环，\(a, b\) 是 A 的两个双边理想，满足 \(a + b = A\) 且 \(a \cap b = \{0\}\) 。试证拓扑环 A 同构于拓扑环 a 与 b 的乘积。
\item
  设 A 是无单位元的拓扑环，B 是由 A 添加一个单位元而得到的环。试证：B 上由 A 的拓扑与 Z 上离散拓扑的乘积所成的拓扑与 B 的环结构相容，并且在 A（B 的一个双边理想）上诱导出所给的拓扑。
\item
  设 \(\mathbf{A}\) 是有单位元的拓扑环。
\end{enumerate}

\begin{enumerate}
\def\labelenumi{\alph{enumi})}
\item
  A 的可逆元所成的集 A* 在 A 中是开的，当且仅当存在 1 的一个邻域 V，它的所有元素都是 A 的可逆元。
\item
  若 A* 在 A 中是开的，则 A 的每个极大（左或右）理想都是闭的，并且 A 的根基是闭的。若 A 除 \(\{o\}\) 与 A 外没有其他闭左理想，则 A 是一个除环。
\item
  设 A 是拓扑域 Q 中由所有有理数 \(k/2^{n}\) （ \(k \in Z, n \in N\) ）组成的子环。试证 A 不含 \(\{0\}\) 与 A 以外的闭理想，但 A 不是域。
\end{enumerate}

\begin{enumerate}
\def\labelenumi{\arabic{enumi})}
\setcounter{enumi}{8}
\tightlist
\item
  \begin{enumerate}
  \def\labelenumii{\alph{enumii})}
  \tightlist
  \item
    元素 \(x\) （相应地，理想 \(\mathfrak{a}\) ）称为拓扑环 \(\mathbf{A}\) 中的拓扑幂零元，如果序列 \((x^n)_{n\geqslant 1}\) 收敛于 o {[}相应地，如果由诸 \(\mathfrak{a}^n\) （ \(n\geqslant 1\) ）组成的滤子基收敛于 o）。拓扑幂零理想的每个元素都是拓扑幂零的。
  \end{enumerate}
\end{enumerate}

\begin{enumerate}
\def\labelenumi{\alph{enumi})}
\setcounter{enumi}{1}
\item
  假设 A 有单位元，且 \(A^{*}\) 在 A 中是开的。试证：若 x 是 A 的任一拓扑幂零元素，则 1 - x 是可逆元（注意，当 n 充分大时 \(1 - x^{n}\) 是可逆元）。由此试证：若（左或右）理想 a 的所有元素都是拓扑幂零的，则 a 包含在 A 的根基中。
\item
  假设 A 完备且有单位元，并且存在由 A 的加法群的子群组成的 o 的基本邻域系。试证：若 x 拓扑幂零，则 1 - x 是 A 中的可逆元。
\end{enumerate}

\(d)\) 在 \(c\) ) 的假设下，试证：若 \(y \in A\) 拓扑幂零，则方程 \(x^{2} + x = y\) 有一个根 \(x \in A\) ，并且 \(x\) 拓扑幂零。

¶ 10) a) 设 B 是有单位元的环，E 是自由 B-模 \(\mathbf{B}_{d}^{(I)}\) ，A 是 E 的自同态环 \(\mathcal{L}(\mathbf{E})\) 。对 E 的每个有限子集 F，令 \(V_{F}\) 为所有这样的 \(u \in A\) 组成的集：\(u(x) = 0\) 对所有 \(x \in F\) 成立。试证诸 \(V_{F}\) 构成 A 上某个豪斯多夫拓扑的 o 的基本邻域系，该拓扑与 A 的环结构相容，且对此拓扑 A 全不连通。

\begin{enumerate}
\def\labelenumi{\alph{enumi})}
\setcounter{enumi}{1}
\tightlist
\item
  取 I = N，并取 B 为一个没有零因子但其根基 R 不是 \(\{0\}\) 的交换环。试证 A 的根基不是闭的。{[}设 \((e_{n})\) 是 E 的典范基。一方面注意：每个 \(u \in A\) ，只要它除有限多个指标外都有 \(u(e_{n}) = 0\) ，且 \(u(E) \subset \Re \cdot E\) ，就属于 A 的根基；另一方面考察 A 的如下元素 \(u_{0}\) ：\(u_{0}(e_{n}) = e_{n} + te_{n+1}\) （其中 \(t \neq 0\) 属于 R，对每个 n）{]}。
\end{enumerate}

¶ 11) a) 有单位元的拓扑环称为 Gelfand 环，如果 A* 在 A 中是开的，并且 A 的拓扑在 A* 上诱导的拓扑与 A* 的乘法群结构相容。每个豪斯多夫拓扑除环都是 Gelfand 环{[}参看习题 20 e){]}。

\begin{enumerate}
\def\labelenumi{\alph{enumi})}
\setcounter{enumi}{1}
\tightlist
\item
  试证：若 A 是 Gelfand 环，则每个赋有乘积拓扑的矩阵环 \(\mathbf{M}_{n}(\mathbf{A})\) （乘积拓扑取在 \(A^{n^{2}}\) 上）也是 Gelfand 环。（对 n 用归纳法证明。）
\end{enumerate}

\begin{enumerate}
\def\labelenumi{\arabic{enumi})}
\setcounter{enumi}{11}
\tightlist
\item
  拓扑环的子集 M 称为右有界的（相应地，左有界的），如果对 A 中 o 的每个邻域 U，存在 o 的邻域 V 使 VM⊂U（相应地，MV⊂U）；若 M 既左有界又右有界，则称 M 为有界的。若 A 中存在 o 的有界邻域，则称 A 的拓扑是局部有界的（且称 A 为局部有界拓扑环）。
\end{enumerate}

\begin{enumerate}
\def\labelenumi{\alph{enumi})}
\item
  每个具有由右理想组成的 o 的基本邻域系的拓扑环都是右有界的。试证：在习题 10 a) 中，若 I 无限且 B 是除环，则环 A 是左有界的，但 A 中不存在 o 的右有界邻域。
\item
  若拓扑环 A 是右有界的（相应地，有界的），且具有由 A 的加法群的子群组成的 o 的基本邻域系，则它具有由右（相应地，双边）理想组成的 o 的基本邻域系。
\item
  有界集的每个有限并是有界的；有界集的闭包是有界的；若 M 与 N 有界，则 \(M + N\) 与 MN 也有界。
\item
  拓扑环 \(\mathbf{A}\) 中的每个准紧集是有界的。
\item
  若 M 与 N 是 A 的两个有界子集，试证映射 \((x, y) \to xy\) （把 \(M \times N\) 的元素映入 A）是一致连续的。
\item
  设 A 是有单位元的豪斯多夫拓扑环。试证：若 \((x_{n})\) 是 A 中趋于 o 的可逆元序列，则诸元素 \(x_{n}^{-1}\) 所成之集既非左有界也非右有界。
\item
  若 A 是有单位元的有界环，试证映射 \(x \rightarrow x^{-1}\) 在 A 的可逆元集 A* 上是一致连续的。h) 试证：若 A 是有单位元的完备豪斯多夫有界环，则可逆元集 A* 与 A 的根基都是闭的{[}利用 g){]}。
\item
  拓扑环的乘积 \(\prod_{t} A_{t}\) 的子集是有界的，当且仅当它的所有投影都是有界的。
\item
  局部有界环的每个子环以及局部有界环的每个乘积都是局部有界环。
\item
  局部有界豪斯多夫环的完备化是局部有界的。
\end{enumerate}

\begin{enumerate}
\def\labelenumi{\arabic{enumi})}
\setcounter{enumi}{12}
\tightlist
\item
  设 A 是有界拓扑环（习题 12），它有单位元，并且 A 的可逆元集 A* 在 A 中是开的。试证 A 的根基是开的。因此，若 A 无根基，则它是离散的；特别地，有界豪斯多夫拓扑除环是离散的。无根基、有单位元且 A* 为开的紧环 A 是有限的；特别地，紧拓扑除环是有限的。
\end{enumerate}

¶ 14) 设 A 是全不连通（*）、无根基且有单位元的紧环。则 A 同构于一个有限单环族的乘积。{[}首先利用 § 6 习题 12 b) 与 § 4 命题 14 证明：A 具有由双边理想组成的 o 的基本邻域系。然后证明存在 A 的双边理想的一个极大集 \(\Phi\) ，使得：(i) 每个双边理想 \(\mathfrak{M} \in \Phi\) 都是 A 的极大开双边理想；(ii) 没有理想 \(\mathfrak{M} \in \Phi\) 包含 \(\Phi\) 中异于 \(\mathfrak{M}\) 的理想的有限交。再证明 A 同构于诸环 \(A/\mathfrak{M}\) （ \(\mathfrak{M}\) 取遍 \(\Phi\) ）的乘积；先利用习题 9 b) 证明所有 \(\mathfrak{M} \in \Phi\) 的交仅由 A 的零元素组成{]}。
(*) 可以证明，有单位元的紧环必然全不连通。

\begin{enumerate}
\def\labelenumi{\arabic{enumi})}
\setcounter{enumi}{14}
\item
  设 A 是全不连通且有单位元的紧环。试证 A 的根基 \(\Re\) 是拓扑幂零的{[}习题 9 a)；利用 § 6 习题 12 b) 与 § 4 命题 14{]}。由此试证 \(\Re\) 是由所有这样的 \(x \in A\) 组成的集：\(ax\) 对每个 \(a \in A\) 都拓扑幂零{[}参看习题 9 c){]}。
\item
  试证：在环 A（相应地，除环 K）上，\(\mathcal{C}\) （与 A 的环结构（相应地，K 的除环结构）相容的一族拓扑 \((\mathcal{T}_i)\) 的上确界）与该结构相容。对拓扑 \(\mathcal{T}_i\) 中的每一个都左（相应地，右）有界的集 M，对 \(\mathcal{C}\) 也左（相应地，右）有界。
\item
  若 K 是非离散的豪斯多夫拓扑除环，则 K 上的两个乘法一致结构没有一个能与 K 的加法一致结构在 \(K^{*}\) 上诱导的一致结构相比较（参看习题 13）。
\item
  设 K 是豪斯多夫拓扑除环，A 是 K 的闭子集，B 是 K 的紧子集且 \(o \notin B\) 。试证 AB 与 BA 在 K 中是闭的（参看 § 4 第 1 目命题 1 推论 1）。* 在实数域 R 中，给出一个例子，其中 A 是闭的，B 是紧的，\(o \in B\) ，而 AB 不是闭的。*
\item
  \begin{enumerate}
  \def\labelenumii{\alph{enumii})}
  \tightlist
  \item
    设 K 是赋有与 K 的环结构相容的豪斯多夫拓扑 \(\mathcal{G}\) 的除环。试证：K 中存在 o 的邻域 V，使 \((\mathbb{CV}) \cup (\mathbb{CV})^{-1} = \mathbf{K}^*\) {[}这等价于 \(\mathbf{V} \cap (\mathbf{V} \cap \mathbf{K}^*)^{-1} = \emptyset\) {]}。
  \end{enumerate}
\end{enumerate}

\begin{enumerate}
\def\labelenumi{\alph{enumi})}
\setcounter{enumi}{1}
\tightlist
\item
  再假设 \(\mathfrak{C}\) 不是离散的。试证：若 \(\mathbf{U}\) 是 \(o\) 在 \(\mathbf{K}\) 中的任一邻域，且 \(\mathbf{U}^*\) 表示 \(\mathbf{U} \cap \mathbf{K}^*\) ，则
\end{enumerate}

\[
\mathbf {U} ^ {*} (\mathbf {U} ^ {*}) ^ {- 1} = \mathbf {K} ^ {*}.
\]

¶ 20) a) 设 K 是赋有与 K 的环结构相容的非离散豪斯多夫拓扑 \(\mathfrak{C}\) 的除环。K 的子集 M 是右（相应地，左）有界的，当且仅当对 o 的每个邻域 U，存在 \(a \in \mathbf{K}^*\) 使 \(a\mathbf{M} \subset \mathbf{U}\) （相应地，Ma \(\subset\) U）；此时，M 对每个豪斯多夫拓扑 \(\mathfrak{C}'\) （它比 \(\mathfrak{C}\) 粗且与 K 的环结构相容）都是右（相应地，左）有界的。

\begin{enumerate}
\def\labelenumi{\alph{enumi})}
\setcounter{enumi}{1}
\item
  若 \(\mathfrak{G}\) 是局部有界的（习题 12），且 U 是关于 \(\mathfrak{G}\) 的 o 的有界邻域，则诸集 \(x\mathrm{U}\) （相应地，\(\mathrm{U}x\) ）构成 o 的（关于 \(\mathfrak{G}\) 的）基本邻域系，这里 x 取遍 \(\mathbf{K}^*\) 。K 的每点具有关于 \(\mathfrak{G}\) 的可数基本邻域系，当且仅当存在以 o 为聚点的 K 中的点列 \((a_n)\) 。
\item
  K 上局部有界拓扑（与 K 的环结构相容的）的任一有限族的上确界是局部有界的。反之，设 \(\mathcal{G}\) 是 K 上局部有界拓扑的某个族 \((\mathfrak{T}_i)\) 的上确界；若它局部有界，则 \(\mathcal{G}\) 是 \((\mathfrak{T}_i)\) 的某个有限子族的上确界。（注意：若 U 是关于 \(\mathcal{G}\) 有界的 o 的邻域，则存在有限的指标组 \((\iota_k)\) ，并且对每个 \(k\) 存在 o 的邻域 \(\mathrm{V}_k\) （它关于 \(\mathcal{G}_{\iota_k}\) 有界），使 \(\bigcap \mathrm{V}_k \subset \mathrm{U}\) ；另一方面，存在 \(a_k \in \mathrm{K}^*\) 使 \(\mathrm{U} \subset a_k \mathrm{V}_{k'}\) 。）
\item
  除环 K 的子集 F 称为拟环，如果：(i) \(F \neq K\) ，\(o \in F\) ，\(1 \in F\) ，\(-F = F\) ，\(FF \subset F\) ；(ii) 存在元素 \(c \in F^{*} = F \cap K^{*}\) 使 \(c(F + F) \subset F\) ；并且 (iii) 对每个 \(x \in F\) ，存在 \(y \in F\) 使 \(yF \subset Fx\) 且 \(Fy \subset xF\) 。试证：若 G 是局部有界且与 K 的环结构相容的非离散豪斯多夫拓扑，则 o 的每个满足 \(UU \subset U\) 且 \(1 \in U\) 的（关于 G 的）对称有界邻域 U 都是拟环，并且 \(U^{*}(U^{*})^{-1} = K^{*}\) {[}习题 19 b){]}。若 V 是关于 G 的 o 的任一对称有界邻域，则由诸元素 \(x \in K\) （它们满足 \(xV \subset V\) ，或 \(Vx \subset V\) ）组成的集 U 是关于 G 的 o 的对称有界邻域，它满足 \(UU \subset U\) 且 \(1 \in U\) ，因而是拟环。
\item
  反之，设 F 是 K 中满足 \(\mathrm{F}^{*}(\mathrm{F}^{*})^{-1} = \mathrm{K}^{*}\) 的拟环。试证：存在非离散且与 K 的环结构相容的豪斯多夫拓扑 G，使 F 是关于 G 的 o 的有界邻域。特别地，可取 F 为 K 的任一子环 A，使 K 是 A 的右分式除环。特别地，若 K = Q 且 F = Z，则 K 上相应的拓扑 G 与 K 的域结构不相容。
\end{enumerate}

¶ 21) a) 设 \(\mathfrak{C}\) 是除环 K 上的局部有界豪斯多夫拓扑，对此拓扑 K 是连通的。试证 K 中存在非零的拓扑幂零元素{[}习题 9 a){]}。（若 U 是 o 的对称有界邻域，t 是 K 中使 \(\mathrm{U}t + t\mathrm{U} \subset \mathrm{U}\) 的非零元素，试利用 § 2 第 2 目命题 6 证明 K 是诸集 \(\mathrm{U}t^{-n}\) 的并。）

\begin{enumerate}
\def\labelenumi{\alph{enumi})}
\setcounter{enumi}{1}
\tightlist
\item
  设 \(\mathfrak{G}\) 是 K 上的局部有界豪斯多夫拓扑，对此拓扑 K 有非零的拓扑幂零元素。试证：存在 o 的邻域 U，使得对 o 的每个邻域 V，存在整数 \(n_0\) ，使得 \(\mathrm{U}^n\subset \mathrm{V}\) 对所有 \(n\geqslant n_0\) 成立；这蕴含 U 的所有元素都拓扑幂零，并且 K 的每点具有关于 \(\mathfrak{G}\) 的可数基本邻域系。（设 \(t\neq 0\) 拓扑幂零，W 是 o 的有界邻域；取 U 为满足 \(\mathrm{UW}\subset \mathrm{Wt}\) 的 o 的有界邻域。）因此，K 的拓扑幂零元素所成之集 T 是关于 \(\mathfrak{G}\) 的 o 的邻域。
\end{enumerate}

¶ 22) 在赋有与 K 的环结构相容的豪斯多夫拓扑 \(\mathfrak{G}\) 的除环 K 中，包含 o 的集 R 称为逆有界的，如果 \((\mathbb{C}R)^{-1}\) 是有界的。若 K 中存在关于 \(\mathfrak{G}\) 的由 o 的逆有界邻域组成的基本系，则称 \(\mathfrak{G}\) 是局部逆有界的。

\begin{enumerate}
\def\labelenumi{\alph{enumi})}
\item
  局部逆有界的拓扑 T 是局部有界的，并且是所有与 K 的环结构相容的豪斯多夫拓扑所成之集中的一个极小元素；此外，C 与 K 的除环结构相容{[}利用习题 19 a){]}。
\item
  在本习题中从此假设 \(\mathcal{G}\) 是 K 上的局部逆有界拓扑。试证：对 K 中 o 的每个邻域 V，\(x \to x^{-1}\) 在 \(\mathbb{C}\mathrm{V}\) 上是一致连续的。由此试证 K 的完备化 \(\hat{\mathbf{K}}\) 是一个拓扑除环，其拓扑是局部逆有界的。
\item
  K 中 o 的每个邻域（关于 \(\mathfrak{G}\) 的）都包含 o 的一个邻域 V，使 \(\mathbf{E}(\mathbf{V}) = (\mathbb{C}\mathbf{V}) \cap (\mathbb{C}\mathbf{V})^{-1}\) 非空。试证：由 K 的加法一致结构与两个乘法一致结构在 \(\mathbf{E}(\mathbf{V})\) 上诱导的三个一致结构重合。由此推出：若 K 完备，则 \(\mathbf{K}^*\) 关于每个乘法一致结构都是完备的。
\item
  试证：若 \(x\) 是 \(\mathbf{K}^*\) 的元素，则下列三个命题中恰有一个为真：(i) \(x\) 拓扑幂零；(ii) \(x^{-1}\) 拓扑幂零；(iii) 序列 \((x^n)_{n\in \mathbb{Z}}\) 有界。{[}通过考察某个有界邻域 \(U\) ，它包含 \(o\) 且满足 \(UU\subset U\) （习题 20 d)），证明：若 o 是序列 \((x^n)_{n\geqslant 0}\) 的聚点，则 \(x\) 拓扑幂零。{]}若 K 含非零的拓扑幂零元素，试证：由 o 与所有这样的 \(x\in \mathbf{K}^*\) （ \(x^{-1}\) 不拓扑幂零）组成的集 B，是 K 中 o 的有界邻域，并且在所有内自同构下不变{[}利用习题 21 b){]}；于是 B 包含所有满足 \(\mathrm{UU}\subset \mathrm{U}\) 的 o 的有界邻域 U；这类在所有内自同构下不变的邻域所成之集有一个最大元素，当 \(\mathbf{K}^*\) 的换位子群有界时，特别地当 K 是域时，它与 B 重合。
\item
  若 U 是 K 中 o 的有界对称邻域且 \(UU \subset U\) ，则存在 \(b \in U^{*}\) 使 \(\mathbf{K}^{*} = \mathbf{U}^{*} \cup ((\mathbf{U}^{*})^{-1} b)\) 。设 \(a \in U^{*}\) 满足 \(Ua \subset bU\) ，并假设序列 \((a^{-n})_{n \geqslant 0}\) 有界。试证：由所有这样的 \(x \in K\) 组成的集 V：xU 包含于诸集 \(Ua^{-n}\) 的并之中，是 o 的有界邻域，满足 \(VV \subset V\) 且 \(\mathbf{K}^{*} = \mathbf{V}^{*} \cup (\mathbf{V}^{*})^{-1}\) 。
\item
  假设 K 不含非零的拓扑幂零元素。试证：K 中存在 o 的有界邻域，它是 K 的子环 A，且使 \(\mathbf{K}^{*} = \mathbf{A}^{*} \cup (\mathbf{A}^{*})^{-1}\) 。{[}从满足 \(\mathbf{K}^{*} = \mathbf{V}^{*} \cup (\mathbf{V}^{*})^{-1}\) 的 o 的有界邻域 V 出发（见 e){]}；若 \(c \in V^{*}\) 满足 \(c(V + V) \subset V\) ，则注意由 V 生成的环 A 包含于诸集 \(Vc^{-n}\) （ \(n \geqslant 0\) ）的并之中。
\end{enumerate}

¶ 23) 若 p 是任一素数，x 是任一非零有理数，令 \(v_{p}(x)\) 表示 p 在 x 的素因子分解中的指数。\(v_{p}\) 从 \(Q^{*}\) 到 Z 中的映射称为 Q 上的 p 进赋值。由所有这样的 \(x \in Q\) 组成的集：x = 0 或 \(v_{p}(x) \geqslant m\) （对某个 \(m \in Z\) ），是 Q 的分式理想 ( \(p^{m}\) )。

\begin{enumerate}
\def\labelenumi{\alph{enumi})}
\tightlist
\item
  试证：诸理想 \((p^{m})\) （ \(m \in Z\) ）构成 Q 中 o 关于某个局部逆有界拓扑（习题 22）\(C_{p}\) 的基本邻域系，该拓扑称为 p 进拓扑。Q 在此拓扑下的完备化 \(Q_{p}\) 是一个域，称为 p 进数域，其元素称为 p 进数。Z 的闭包 \(Z_{p}\) （在 \(Q_{p}\) 中）是 \(Q_{p}\) 中的一个紧开主理想整环，其中 \(p = pZ_{p}\) 是唯一的非零素理想；\(Z_{p}/p\) 同构于素域 \(F_{p} = Z/(p)\) ，并且更一般地，加法商群 \(p^{m}/p^{n}\) 同构于
\end{enumerate}

\[
\mathbf {Z} / (p ^ {n - m}) \text {   if   } m <   n.
\]

\(\mathbf{Z}_p\) 的元素称为 \(p\) 进整数。

\begin{enumerate}
\def\labelenumi{\alph{enumi})}
\setcounter{enumi}{1}
\item
  试证：加法群 \(Q_{p}\) 到其自身的每个连续同态都具有 \(x \to ax\) 的形式，其中 \(a \in Q_{p}\) {[}若 f 是这样一个同态，试证 \(f(rx) = rf(x)\) 对所有 \(r \in Q\) 成立，然后在 \(Q_{p}\) 中取极限{]}。
\item
  试证：每个紧子群 \(G \neq \{o\}\) （ \(Q_{p}\) 的加法群的子群）都与某个 \(p^{n} (n \in Z)\) 相同，并且加法群 \(Q_{p}\) 除 \(Q_{p}\) 自身外没有非紧的闭子群。{[}设 m 是使 \(G \subset p^{m}\) 的最大整数；通过考察商群 \((G + p^{n})/p^{n}\) （当 n \textgreater{} m 时），证明 \(G + p^{n} = p^{m}\) ，然后利用 § 3 第 1 目的公式 (1){]}。
\end{enumerate}

¶ 24) a) 试证：乘法群 \(\mathbf{Q}_p^*\) （ \(p\) 进数域 \(\mathbf{Q}_p\) 的乘法群）的每个子群都同构于乘法群 \(\mathbf{U}\) （ \(\mathbf{Z}_p\) 的可逆元的乘法群）的一个子群与一个同构于 \(\mathbf{Z}\) 或 \(\{o\}\) 的离散加法子群的乘积。

\begin{enumerate}
\def\labelenumi{\alph{enumi})}
\setcounter{enumi}{1}
\item
  试证：U 的子群 \(V = 1 + p\) 的紧子群就是诸群 \(1 + p^{n}\) {[}推理与习题 23 c) 的相同{]}。
\item
  试证：对每个 \(a \in \mathbf{U}\) ，序列 \((a^{p^n})_{n \in \mathbb{N}}\) 趋于一个极限 \(\alpha \equiv a \pmod{\mathfrak{p}}\) ，并且 \(\alpha^p = \alpha\) {[}证明 \(a^{p^n} \equiv a^{p^{n-1}} \pmod{\mathfrak{p}^n}\) （对 \(n\) 用归纳法）{]}。试证：多项式 \(\mathbf{X}^{p-1} - 1\) （在 \(\mathbf{Q}_p\) 的一个代数闭扩张中）的所有根都属于 \(\mathbf{Q}_p\) ，并且模 \(\mathfrak{p}\) 两两不同余{[}把前面所述应用于同余式 \(x^{p-1} - 1 \equiv 0 \pmod{p}\) 的根{]}。若 \(p > 2\) 且 \(d\) 是 \(n\) 与 \(p - 1\) 的最大公因子，则多项式 \(\mathbf{X}^n - 1\) 恰有 \(d\) 个根（在 \(\mathbf{Q}_p\) 中），它们都是 \(\mathbf{X}^d - 1\) 的根（同样的方法：取 \(a\) 为 \(\mathbf{X}^n - 1\) 的一个根）。
\end{enumerate}

由此试证：若 \(p > 2\) ，则乘法群 \(\mathbf{Q}_p^*\) 的每个紧子群都是 \(\mathbf{U}\) 的一个有限子群{[}由 \((p - 1)\) 次单位根组成{]}与一个形如 \(1 + \mathfrak{p}^n\) 的子群的直积{[}利用 \(b\){]}。当 \(p = 2\) 时这些结果应如何修改？（让群 \(1 + \mathfrak{p}^2\) 起先前由 \(1 + \mathfrak{p}\) 所起的作用。）

\begin{enumerate}
\def\labelenumi{(\arabic{enumi})}
\setcounter{enumi}{24}
\tightlist
\item
  \begin{enumerate}
  \def\labelenumii{\alph{enumii})}
  \tightlist
  \item
    设 \(a \neq o\) 是一个 \(p\) 进数。则映射 \(n \to a^n\) 在 \(\mathbf{Z}\) 上（作为 \(\mathbf{Q}_p\) 的子空间）连续，当且仅当 \(a \in 1 + \mathfrak{p}\) 。若此条件满足，试证 \(n \to a^n\) 在 \(\mathbf{Z}\) 上一致连续，并且可延拓为 \(\mathbf{Z}_p\) 到 U 中的连续同态，记作 \(x \to a^x\) ，当 \(a \neq 1\) 时它是单射。若 \(p > 2\) 且 \(m\) 是使 \(a \in 1 + \mathfrak{p}^m\) 的最大整数，则 \(x \to a^x\) 是 \(\mathbf{Z}_p\) 到 \(1 + \mathfrak{p}^m\) 上的一个同构。当 \(p = 2\) 时这个命题应如何修改？b) 试证：\(\mathbf{Z}_p\) 到 U 中的每个连续同态都具有 \(x \to a^x\) 的形式，其中 \(a \in 1 + \mathfrak{p}\) {[}若 \(f\) 是这样一个同态且 \(f(1) = a\) ，则 \(f(n) = a^n\) 对 \(n \in \mathbf{Z}]\) 成立。
  \end{enumerate}
\end{enumerate}

\begin{enumerate}
\def\labelenumi{\alph{enumi})}
\setcounter{enumi}{2}
\item
  试证：映射 \((x, y) \to x^y\) 在 \((1 + \mathfrak{p}) \times \mathbf{Z}_p\) 上连续。若 \(p > 2\) 、\(b \in \mathbf{Z}_p\) ，且 \(m\) 是使 \(b \in \mathfrak{p}^m\) 的最大整数，则映射 \(x \to x^b\) 是乘法群 \(1 + \mathfrak{p}\) 到乘法子群 \(1 + \mathfrak{p}^{m+1}\) 上的一个同构{[}利用习题 24 b){]}。当 \(p = 2\) 时这个命题应如何修改？
\item
  若 \(n\) 是与 \(p - 1\) 和 \(p\) 都互素的整数，试证映射 \(x \to x^n\) 是乘法群 U 的一个自同构{[}利用习题 24 c){]}。
\end{enumerate}

\begin{enumerate}
\def\labelenumi{\arabic{enumi})}
\setcounter{enumi}{25}
\tightlist
\item
  \begin{enumerate}
  \def\labelenumii{\alph{enumii})}
  \tightlist
  \item
    设 \(p, q\) 是两个不同的素数，\(\mathfrak{C}\) 是 \(\mathbf{Q}\) 上作为拓扑 \(\mathfrak{C}_p\) 与 \(\mathfrak{C}_q\) 的上确界的拓扑。则 \(\mathfrak{C}\) 是局部有界的{[}习题 20 c){]}，并且与 \(\mathbf{Q}\) 的域结构相容。关于此拓扑，两个序列 \(((p/q)^n)_{n\geqslant 0}, ((q/p)^n)_{n\geqslant 0}\) 中没有一个是有界的；序列 \((1/p^n)_{n\geqslant 0}\) 不是有界的，但有界的序列 \((p^n)_{n\geqslant 0}\) 并不趋于 o{[}参看习题 22 d){]}。试证 \(\mathbf{Q}\) 关于拓扑 \(\mathfrak{C}\) 的完备化同构于拓扑域 \(\mathbf{Q}_p\) 与 \(\mathbf{Q}_q\) 的乘积。
  \end{enumerate}
\end{enumerate}

\begin{enumerate}
\def\labelenumi{\alph{enumi})}
\setcounter{enumi}{1}
\tightlist
\item
  若 P 是所有素数组成的集，则拓扑 \(G_{0}\) （它是诸拓扑 \(G_{p}\) （ \(p \in P\) ）的上确界）与 Q 的域结构相容，但不是局部有界的{[}习题 20 c){]}。Q 关于拓扑 \(G_{0}\) 的完备化是什么？
\end{enumerate}

